\paragraph{On-Policy Distillation.} Distillation for LLMs is broadly categorized into off-policy and on-policy settings. While off-policy distillation (e.g., SFT) suffers from exposure bias due to the mismatch between static teacher-generated data and student-generated rollouts at inference \citep{agarwal2024policy}, OPD mitigates this issue by training on trajectories sampled directly from the student's current policy under a reverse KL objective \citep{gu2024minillm}. Recently, OPD has since appeared in some reasoning and post-training recipes, such as Qwen3 \citep{yang2025qwen3}, Deepseek-V4 \citep{xu2026deepseek} and GLM-5 \citep{zeng2026glm}. Building on these successes, a growing body of work has investigated various algorithmic variants of OPD \citep{oh2026kl, yang2026prune}, optimized its practical training recipes \citep{li2026filter, ko2026scaling}, and integrated it into large-scale post-training pipelines \citep{ma2026mopd}. Moreover, several studies have further analyzed the underlying mechanisms of OPD \citep{wang2026not, li2026rethinking}. 

\paragraph{Data Selection for LLM Post-Training.} Selecting a small yet effective subset of training data has become an important problem for improving the efficiency of LLM post-training, with most efforts focusing on data selection for SFT \citep{ivison2025large} and RLVR \citep{wu2026single}. Existing approaches include LLM-based quality assessment \citep{chen2024alpagasus}, gradient-based selection \citep{xia2024less} and more. However, data selection for OPD remains underexplored. Our work is inspired by recent advances in RLVR, which show that RLVR can achieve significant reasoning gains with as few as a single training example \citep{wang2026reinforcement}. We extend this exploration to the OPD paradigm, discovering different and interesting phenomena. A concurrent work by \citet{fu2026rethinking} also investigates 1-shot OPD and analyzes its efficacy from a state-space coverage perspective. Differently, we analyze 1-shot OPD from the perspective of CoT trajectory and reasoning patterns. Then we propose a simple and effective data selection method that selects only hard examples for training, where even ``unsolvable'' examples completely exceeding the teacher's capability can be successfully used.